%%________________________________________________________________________
%% LEIM | PROJETO
%% 2022 / 2013 / 2012
%% Modelo para relatório
%% v04: alteração ADEETC para DEETC; outros ajustes
%% v03: correção de gralhas
%% v02: inclui anexo sobre utilização sistema controlo de versões
%% v01: original
%% PTS / MAR.2022 / MAI.2013 / 23.MAI.2012 (construído)
%%________________________________________________________________________




%%________________________________________________________________________
\myPrefaceChapter{Resumo}
%%________________________________________________________________________

O projeto Deep-Audio surge da necessidade de transcrição de áudios no trabalho da Polícia Judiciária, um processo demorado e cansativo quando feito manualmente. 

Desenvolvemos uma solução eficiente que facilitasse esta tarefa, utilizando modelos de inteligência artificial \textit{open-source}, como o Whisper da OpenAI, para converter áudio em texto, e o Pyannote, para identificar os oradores presentes no áudio. Além disso, o sistema permite corrigir possíveis falhas do modelo e gerar datasets para futuro treino de modelos.

As ideias mais relevantes incluem a criação de uma aplicação \textit{web} intuitiva, que transcreve áudios importados, permite a organização dos áudios em pastas, exibe as transcrições em formato de conversa (\textit{chat}) e permite edições quer do texto, quer do nome dos oradores por parte do utilizador.

O desenvolvimento do Deep-Audio seguiu uma abordagem iterativa e incremental, dividida em fases claras de forma a garantir uma melhor organização do fluxo de trabalho ao longo do tempo. 

Explorámos as tecnologias utilizadas para a transcrição e para a identificação de oradores, através da criação de protótipos para testar a integração entre as várias partes do sistema. Foram realizados testes incrementais em cada componente, desde o \textit{upload} de áudios até a exportação da transcrição final para um documento.

A interface gráfica foi desenvolvida com foco na usabilidade, incorporando \textit{feedback} dos utilizadores para ajustes nas funcionalidades e na própria interface. 

Esta abordagem metodológica permitiu um desenvolvimento estruturado, que possibilitou alcançar as metas pretendidas para o projeto.





%%________________________________________________________________________
\myPrefaceChapter{Abstract}
%%________________________________________________________________________

The Deep-Audio project arose from the need to transcribe audio in the work of the Judicial Police, a time-consuming and tiring process when done manually.

We developed an efficient solution that would make this task easier, using \textit{open-source} artificial intelligence models, such as Whisper from OpenAI, to convert audio into text, and Pyannote, to identify the speakers present in the audio. In addition, the system makes it possible to correct possible flaws in the model and generate datasets for future model training.

The most important ideas include the creation of an intuitive \textit{web} application, which transcribes imported audio, organizes the audio in folders, displays the transcripts in conversation format (\textit{chat}) and allows the user to edit both the text and the names of the speakers.

The development of Deep-Audio followed an iterative and incremental approach, divided into clear phases to ensure better organization of the workflow over time.

We explored the technologies used for transcription and speaker identification by creating prototypes to test the integration between the various parts of the system. Incremental tests were carried out on each component, from the \textit{upload} of audios to the export of the final transcript to a document.

The graphic interface was developed with a focus on usability, incorporating feedback from users to adjust the functionalities and the interface itself.

This methodological approach enabled structured development, which made it possible to achieve the project's goals.